\section*{अध्याय \ref{ch:b1:radiation-heat-transport} --- विकिरण और ऊष्मा का परिवहन}
% values: calc/ch27.py

\begin{solution}{exo:b1:radiation-heat-transport:1}
(क)~$S_0/\pi = 1361/\pi = \qty{433.2}{W.m^{-2}}$। (ख)~विषुव पर भूमध्य रेखा पर
$\bar Q = (S_0/\pi)(a/r)^2$, इसलिए $(a/r)^2 = 436.7/433.2 = 1.0080$ और $r/a =
1.0080^{-1/2} = 0.9960$: पृथ्वी अपनी माध्य दूरी से \qty{0.4}{\percent} निकट सूर्य के है। (ग)~$\bar Q = 433.2 \times 0.9683 \times \cos\ang{23.44} =
\qty{384.9}{W.m^{-2}}$।
\end{solution}

\begin{solution}{exo:b1:radiation-heat-transport:2}
(क)~$415.6/172.4 = 2.41$। (ख)~समांतर रेखाओं $\pm\varphi$ के बीच की पट्टी का क्षेत्रफल
$2\pi a_\oplus^2 \times 2\sin\varphi$ है, जो गोले का $\sin\varphi$ भाग है: $\sin\ang{37.0} = 0.602$, \qty{60}{\percent}। (ग)~\omterm{def:b1:spherical-astronomy:geographic}{अक्षांश} के अक्ष पर ध्रुवों के निकट की पट्टियाँ भूमध्य रेखा के निकट की पट्टियों जितनी चौड़ी दिखती हैं, पर वे बहुत कम क्षेत्रफल ढकती हैं: \ang{60} के ध्रुव की ओर का क्षेत्र केवल $1 - \sin\ang{60} =
\qty{13.4}{\percent}$ सतह का है। $\sin\varphi$ के विरुद्ध, उष्णकटिबंध अक्ष का अधिकांश भरते हैं।
\end{solution}

\begin{solution}{exo:b1:radiation-heat-transport:3}
(क)~$T_e = (\text{OLR}/\sigma)^{1/4}$: \qty{256.4}{K}, \qty{241.3}{K} और
\qty{254.5}{K}। (ख)~शोषित सूर्य-प्रकाश: $287.7/75.8 = 3.80$; OLR: $244.9/192.3 =
1.27$। (ग)~\omterm{def:b1:radiation-heat-transport:net-radiation}{शुद्ध विकिरण} शोषित सूर्य-प्रकाश का अनुसरण करता है: उष्णकटिबंधों में धनात्मक और बड़ा, उच्च अक्षांशों पर तीव्र ऋणात्मक। OLR, जो थोड़ा बदलता है, शोषित सूर्य-प्रकाश के गिरने की भरपाई नहीं करता।
\end{solution}

\begin{solution}{exo:b1:radiation-heat-transport:4}
(क)~$2.82 \times 365.25 = \qty{1030}{mm}$, प्रति वर्ष \qty{1.03}{m}। (ख)~$1.030 \times
78.1 = \qty{80.4}{W.m^{-2}}$।
(ग)~पुनर्विश्लेषण का $LE$ \qty{2}{W.m^{-2}} बड़ा है। कई वर्षों में वायुमंडल जल-वाष्प जमा करता नहीं रह सकता: वाष्पित सारा जल नीचे लौट आता है, इसलिए वैश्विक माध्य में $LE$ और $LP$ बराबर होने चाहिए। अंतर मॉडल का एक असंतुलन है, उसके विकिरणीय असंतुलन की तरह।
\end{solution}

\begin{solution}{exo:b1:radiation-heat-transport:5}
(क)~$\bar Q = S_0(a/r)^2\sin\delta = 1361 \times 0.9683 \times 0.3978 =
\qty{524.2}{W.m^{-2}}$। (ख)~$1361 \times 1.0334 \times 0.3978 =
\qty{559.5}{W.m^{-2}}$; अनुपात $1.0334/0.9683 = 1.067$ है। (ग)~ध्रुव पर सूर्य-प्रकाश 24 में से 24 घंटे स्थिर ऊँचाई पर रहता है, इसलिए दिन पर उसका माध्य $\sin h_\odot$
$\sin\delta = 0.398$ है। भूमध्य रेखा पर सूर्य दोपहर को ऊँचा होता है पर आधे समय नीचे रहता है; 24 घंटों पर $\sin h_\odot$ का माध्य $\cos\delta/\pi =
0.292$ है। अनुपात $\pi\tan\delta = 1.362$ है
(\cref{prop:b1:radiation-heat-transport:polar-summer})।
\end{solution}

\begin{solution}{exo:b1:radiation-heat-transport:6}
(क)~$\cos h_0 = -\tan\ang{40}\tan\ang{23.44} = -0.3638$, $h_0 = \ang{111.33} =
\qty{1.943}{rad}$; दिन $24 \times 111.33/180 = \qty{14.84}{h}$ तक रहता है; $\bar Q =
(1361/\pi) \times 0.9683 \times (1.943 \times 0.6428 \times 0.3978 + 0.7660 \times
0.9175 \times 0.9315) = \qty{483.0}{W.m^{-2}}$। (ख)~$\cos h_0 = +0.3638$, $h_0 =
\ang{68.67}$, \qty{9.16}{h}; $\bar Q = \qty{155.9}{W.m^{-2}}$। (ग)~अनुपात
$483.0/155.9 = 3.10$ है। दूरी $0.9683/1.0334 = 0.937$ योगदान देती है: वह उत्तरी ग्रीष्म के विरुद्ध काम करती है। ज्यामिति, सूर्य की ऊँचाई और दिन की लंबाई, $3.10/0.937 = 3.31$ योगदान देती है, मौसमी विरोधाभास का पूरा भाग और उससे अधिक।
\end{solution}

\begin{solution}{exo:b1:radiation-heat-transport:7}
(क)~$B = H/LE$: भूमध्य रेखा पर \num{0.093}, \ang{25} उ पर \num{0.374},
\ang{45} उ पर \num{0.338}, \ang{70} उ पर \num{-0.152}। (ख)~\ang{25} उ: ज़ोनल
माध्य में महान गरम मरुस्थल हैं, सूखी ज़मीन जिसमें वाष्पित करने को थोड़ा जल है, जहाँ गरम ज़मीन हवा को इसके बजाय गरम करती है। (ग)~$H < 0$ का अर्थ है कि \omterm{def:b1:heat-budget:heat-flow}{ऊष्मा-प्रवाह} हवा से सतह की ओर होता है: वार्षिक माध्य में बर्फ़ और समुद्री बर्फ़ उनके ऊपर की हवा से ठंडी हैं, जो निचले अक्षांशों से लाई गई हवा से गरम होती है।
\end{solution}

\begin{solution}{exo:b1:radiation-heat-transport:8}
(क)~आकलित: $\alpha = 100/340 = 0.294$, $R = 340 - 100 - 239 =
\qty{+1}{W.m^{-2}}$ (\num{0.7} पूर्णांकित)। पुनर्विश्लेषण: $\alpha = 117.0/341.9 =
0.342$, $R = 341.9 - 117.0 - 237.7 = \qty{-12.8}{W.m^{-2}}$ इन पूर्णांकित पदों के साथ,
\qty{-12.9}{W.m^{-2}} अपरिवर्तित पदों के साथ \cref{ex:b1:radiation-heat-transport:global} में। (ख)~पुनर्विश्लेषण $224.8$ को $240$ के विरुद्ध सोखता है, \qty{15.2}{W.m^{-2}} कम, क्योंकि वह \qty{17}{W.m^{-2}} अधिक परावर्तित करता है; उसका OLR केवल \qty{1.3}{W.m^{-2}} कम है। अधिक परावर्तन असंतुलन समझाता है। (ग)~$R$ को दक्षिणी ध्रुव से उत्तरी ध्रुव तक समाकलित करने पर वैश्विक समाकल
$-12.9 \times 4\pi a_\oplus^2 = -12.9 \times
5.10\times10^{14} = \qty{-6.6}{PW}$ मिलता है। उत्तरी ध्रुव के आर-पार परिवहन असंभव है: उसके परे कुछ नहीं है, समांतर रेखा एक बिंदु है। प्रस्ताव तभी सही है यदि $\int R\,\mathrm dA = 0$ हो, जो स्थायी अवस्था माँगती है।
\end{solution}

\begin{solution}{exo:b1:radiation-heat-transport:9}
(क)~क्षेत्रफल $2\pi a_\oplus^2 \times 2\sin\ang{30} = 2\pi a_\oplus^2 =
\qty{2.55e14}{m^2}$, पृथ्वी का आधा; आधिक्य $33 \times 2.55\times10^{14} =
\qty{8.4}{PW}$। (ख)~हर समांतर रेखा के आर-पार \qty{4.2}{PW}। (ग)~दोनों मिलकर \qty{8.5}{PW} जोड़ते हैं, यानी आधिक्य, पर वे असमान हैं: दक्षिणी गोलार्ध उत्तरी से \qty{0.5}{PW} अधिक निर्यात करता है। भूमध्य रेखा पर आधिक्य समान नहीं बाँटा जाता: इस पुनर्विश्लेषण में उत्तरी उष्णकटिबंध के आधिक्य का \qty{0.36}{PW} भूमध्य रेखा को दक्षिण की ओर पार करता है।
\end{solution}

\begin{solution}{exo:b1:radiation-heat-transport:10}
(क)~दैनिक माध्य को $\bar Q = (a/r)^2 F(\varphi, \delta)$ लिखें, जहाँ $F$ का
कोष्ठक \cref{thm:b1:radiation-heat-transport:daily} का गुणा $S_0/\pi$। इस दौरान
$\mathrm dt$ में प्रति वर्ग मीटर मिली ऊर्जा $\bar Q\,\mathrm dt$ है, और
$\mathrm dt = r^2\,\mathrm d\lambda/h$ इसलिए $\bar Q\,\mathrm dt = (a^2/h)
F(\varphi, \delta(\lambda))\,\mathrm d\lambda$: गुणक $r^2$ निरस्त हो जाता है।
(ख)~एक वर्ष पर ऊर्जा $(a^2/h)\int_0^{2\pi} F(\varphi,
\delta(\lambda))\,\mathrm d\lambda$ है। सूत्र अपरिवर्तित रहता है जब $\varphi$ और
$\delta$ एक साथ चिह्न बदलते हैं, $F(-\varphi, \delta) = F(\varphi, -\delta)$, और
$\delta(\lambda + \pi) = -\delta(\lambda)$, क्योंकि $\sin\delta = \sin\varepsilon
\sin\lambda$। समाकल में $\lambda$ को $\pi$ से खिसकाने पर दिखता है कि $-\varphi$
वही ऊर्जा पाता है जो $\varphi$। (ग)~उत्तरी वसंत और ग्रीष्म
\num{7.6} दिन लंबे रहते हैं, क्योंकि पृथ्वी तब अपसौर के निकट होती है और धीरे चलती है; दक्षिणी ग्रीष्म अधिक चमकीली है, पर छोटी।
(क) से, दोनों प्रभाव ठीक एक-दूसरे को निरस्त करते हैं।
\end{solution}

\begin{solution}{exo:b1:radiation-heat-transport:11}
(क)~$\langle (a/r)^2\rangle = \frac1P\int_0^P \frac{a^2}{r^2}\,\mathrm dt =
\frac1P\int_0^{2\pi}\frac{a^2}{r^2}\frac{r^2}{h}\,\mathrm d\lambda = \frac{2\pi
a^2}{hP} = \frac{2\pi a^2}{2\pi a b} = \frac ab$, और $b = a\sqrt{1-e^2}$।
(ख)~तात्कालिक वैश्विक माध्य $\frac14 S_0(a/r)^2$ का माध्य, यानी $\frac14 S_0
\langle(a/r)^2\rangle = S_0/(4\sqrt{1-e^2})$ है। (ग)~अब $340.25/\sqrt{1 - 0.0167^2} =
\qty{340.30}{W.m^{-2}}$, $S_0/4$ से \qty{0.05}{W.m^{-2}} या \qty{0.014}{\percent} ऊपर।
\end{solution}

\begin{solution}{exo:b1:radiation-heat-transport:12}
(क)~वायुमंडल $0.78 \times 4.10 = \qty{3.20}{PW}$, महासागर $0.22 \times 4.10 =
\qty{0.90}{PW}$। (ख)~अकेला वायुमंडल अधिक ले जाता पाया गया, \qty{5.0}{PW},
पुनर्विश्लेषण के वायुमंडल और महासागर दोनों मिलाकर \qty{4.02}{PW} से: पुनर्विश्लेषण का
कुल परिवहन कम से कम एक-पाँचवाँ कम है ($1 - 4.02/5.0 = 0.20$)।
उसका \omterm{def:b1:radiation-heat-transport:net-radiation}{शुद्ध विकिरण} उष्णकटिबंधों और ध्रुवों के बीच बहुत कम बदलता है। (ग)~$R$ से एक स्थिर संख्या घटाने पर उसका वैश्विक समाकल शून्य होता है, पर \omterm{def:b1:spherical-astronomy:geographic}{अक्षांश} के साथ उसका बदलाव अपरिवर्तित रहता है;
त्रुटि उसी बदलाव में है, हर \omterm{def:b1:spherical-astronomy:geographic}{अक्षांश} पर मॉडल के बादलों और सतहों के एल्बिडो में, केवल वैश्विक माध्य में नहीं।
\end{solution}

\begin{solution}{pb:b1:radiation-heat-transport:1}
\begin{enumerate}
  \item जून के लिए: $\cos h_0 = -0.3638$, $h_0 = \ang{111.33}$, $\bar Q =
        \qty{483.0}{W.m^{-2}}$। दिसंबर: $h_0 = \ang{68.67}$, $\bar Q =
        \qty{155.9}{W.m^{-2}}$। मार्च: $h_0 = \ang{90}$, $\bar Q = (1361/\pi)
        \times 1.0081 \times \cos\ang{40} = \qty{334.6}{W.m^{-2}}$।
  \item जून: $\tan\ang{80}\tan\ang{23.44} = 2.46 > 1$, ध्रुवीय दिवस, $h_0 = \pi$:
        $\bar Q = 1361 \times 0.9683 \times \sin\ang{80}\sin\ang{23.44} =
        \qty{516.3}{W.m^{-2}}$। दिसंबर: ध्रुवीय रात, $\bar Q = 0$। मार्च:
        $(1361/\pi) \times 1.0081 \times \cos\ang{80} = \qty{75.8}{W.m^{-2}}$।
  \item \ang{80} उ पर सूर्य नहीं डूबता ($h_0 = \ang{180}$, 24 घंटे), जबकि
        \ang{40} उ पर वह \qty{14.8}{h} चमकता है। उसकी ऊँचाई कम है, \ang{13.4} और \ang{33.4} के बीच, पर पद $h_0\sin\varphi\sin\delta$, जो दिन की लंबाई और \omterm{def:b1:spherical-astronomy:geographic}{अक्षांश} के साथ बढ़ता है, जीतता है।
  \item पट्टी के मान केंद्रीय मानों से \qty{0.5}{\percent}, \qty{0.8}{\percent}, \qty{1.7}{\percent}, \qty{1.7}{\percent} और
        \qty{1.3}{\percent} अधिक हैं। पहला, पुनर्विश्लेषण थोड़ा बड़ा सौर विकिरण उपयोग करता है (प्रश्न~5), \qty{0.5}{\percent} अधिक। दूसरा, पट्टी का माध्य केंद्र का मान नहीं है: हर पट्टी का भूमध्यरेखीय आधा ध्रुवीय आधे से अधिक क्षेत्रफल और अधिक सूर्य-प्रकाश रखता है।
  \item \cref{prop:b1:radiation-heat-transport:global} से, $S_0 = 4\sqrt{1-e^2}
        \times 341.9 = \qty{1367}{W.m^{-2}}$, \qty{0.47}{\percent} अधिक \qty{1361}{W.m^{-2}} से।
  \item $A = 2\pi a_\oplus^2(1 - \sin\ang{40}) = 2\pi \times (6.371\times10^6)^2
        \times 0.3572 = \qty{9.110e13}{m^2}$, पृथ्वी का भाग $(1 - \sin\ang{40})/2 =
        0.179$।
  \item $(1 - \alpha)\bar Q$: \num{198.0}, \num{159.3}, \num{124.7}, \num{92.5},
        \num{74.7} \unit{W.m^{-2}}।
  \item $R$: \num{-28.1}, \num{-51.0}, \num{-76.6}, \num{-104.9}, \num{-116.9}
        \unit{W.m^{-2}}। पाँचों पट्टियाँ घाटे में हैं; ध्रुवीय पट्टी
        \ang{80}--\ang{90} सबसे बड़ी है।
  \item $T_e = (\text{OLR}/\sigma)^{1/4}$: \num{251.3}, \num{246.8}, \num{244.1},
        \num{242.9}, \num{241.1} \unit{K}।
  \item शोषित: $198.0/74.7 = 2.65$; OLR: $226.1/191.6 = 1.18$।
  \item सूर्य नीचा है, और तिरछे कोण पर आने वाला प्रकाश बादलों और पानी से अधिक परावर्तित होता है; बर्फ़ और समुद्री बर्फ़, चमकीली सतहें, ध्रुव की ओर ज़मीन और समुद्र का अधिक भाग ढकती हैं।
  \item OLR $T_e^4$ के अनुपात में चलता है, और $T_e$ टोपी भर में केवल \qty{10}{K}, \qty{4}{\percent} ही बदलता है: उच्च \omterm{def:b1:spherical-astronomy:geographic}{अक्षांश} अपने सूर्य-प्रकाश भर से जितना अनुमत होता उससे कहीं अधिक गरम रखे जाते हैं, दक्षिण से लाई गई ऊष्मा द्वारा।
  \item \unit{m^2} में: \num{3.143e13}, \num{2.550e13}, \num{1.879e13},
        \num{1.151e13}, \num{0.387e13}; योग \qty{9.110e13}{m^2} है।
  \item $R_iA_i$: \num{-0.883}, \num{-1.300}, \num{-1.439}, \num{-1.207},
        \num{-0.453} \unit{PW}; टोपी: \qty{-5.28}{PW}।
  \item $\alpha_\text{टोपी} = \sum\alpha_i\bar Q_iA_i/\sum\bar Q_iA_i = 0.400$,
        भूमंडल के \num{0.342} के विरुद्ध।
  \item $-12.90 \times 5.10\times10^{14} = \qty{-6.58}{PW}$: पृथ्वी उस दर से ऊष्मा खोती और ठंडी होती। आकलित असंतुलन है
        $+0.7 \times 5.10\times10^{14} = \qty{+0.36}{PW}$, अठारह गुना छोटा और विपरीत चिह्न का: पुनर्विश्लेषण का असंतुलन मॉडल की त्रुटि है।
  \item सुधारा गया $R_i + 12.90$: \num{-15.2}, \num{-38.1}, \num{-63.7},
        \num{-92.0}, \num{-104.0} \unit{W.m^{-2}}; शक्तियाँ \num{-0.478},
        \num{-0.972}, \num{-1.197}, \num{-1.059}, \num{-0.403} \unit{PW}; टोपी:
        \qty{-4.108}{PW}, \qty{45.1}{W.m^{-2}} का माध्य घाटा।
  \item टोपी की ऊर्जा वर्ष-दर-वर्ष नहीं बदलती, इसलिए वह जो वायुमंडल के ऊपरी भाग पर खोती है वह उसमें प्रवेश करनी चाहिए, और प्रवेश का एकमात्र रास्ता
        40वीं समांतर रेखा के आर-पार है: $T(\ang{40}) = -2\pi a_\oplus^2\int_{\ang{40}}^{\ang{90}}
        R\cos\varphi\,\mathrm d\varphi$।
  \item $T(\ang{40}\,\text{उ}) = +\qty{4.11}{PW}$, उत्तर की ओर।
  \item $LP = 2.4655\times10^6 \times 1.95/86400 = \qty{55.6}{W.m^{-2}}$, जो
        $LE = \qty{42.4}{W.m^{-2}}$ से अधिक है: टोपी पर उससे अधिक वर्षा गिरती है जितना वह वाष्पित करती है; अंतर वाष्प के रूप में पहुँचता है।
  \item $(LE - LP) - 2.07 = 42.4 - 55.6 - 2.07 = \qty{-15.3}{W.m^{-2}}$; टोपी के क्षेत्रफल से गुणा करने पर: \qty{-1.40}{PW}। वायुमंडल \ang{40} उ के आर-पार \qty{1.40}{PW} गुप्त ऊष्मा लाता है।
  \item $1.40/4.11 = 0.34$: एक तिहाई।
  \item महासागर $0.22 \times 4.11 = \qty{0.90}{PW}$; वायुमंडल $0.78 \times 4.11 =
        \qty{3.20}{PW}$, जिसमें \qty{1.40}{PW} गुप्त और \qty{1.80}{PW}
        वायु की संवेदी ऊष्मा और स्थितिज ऊर्जा के रूप में ले जाया गया।
  \item निचली सीमा। \ang{43} उ पर पुनर्विश्लेषण का कुल, \qty{4.02}{PW}, अध्ययन में अकेले वायुमंडल, \qty{5.0}{PW}, से कम है; सच्चा कुल अधिक है।
  \item लगभग \qty{4.11}{PW}, \qty{4.11e15}{W}, पुनर्विश्लेषण के विकिरण के अनुसार, जिसमें से एक तिहाई गुप्त ऊष्मा है; सत्य मान और भी अधिक है।
\end{enumerate}
\end{solution}
